\chapter{REVUE DE LITTERATURE}
\thispagestyle{empty}

\section{Cadre general du resume automatique}

Le resume automatique est un domaine important du traitement automatique du langage naturel. Son objectif est de produire une version plus courte d'un texte tout en conservant les informations principales. Les premieres approches etaient surtout extractives: elles selectionnaient les phrases les plus importantes du document source et les assemblaient pour former un resume. Ces methodes etaient utiles pour reduire rapidement la taille d'un texte, mais elles produisaient souvent des resumes peu naturels, avec des repetitions ou une coherence limitee entre les phrases. Avec le developpement des modeles neuronaux, le resume abstractive est devenu plus central. Dans cette approche, le modele ne copie pas seulement des phrases du texte source; il peut reformuler l'information et produire de nouvelles phrases plus proches d'un resume humain \cite{Nallapati2016Abstractive,See2017PointerGenerator}.

Les travaux de Nallapati et al. ont montre que les architectures sequence-a-sequence pouvaient etre adaptees au resume abstractive, mais aussi que cette tache restait difficile a cause de la longueur des documents, de la selection de l'information et de la generation d'un texte coherent \cite{Nallapati2016Abstractive}. See et al. ont ensuite propose les reseaux pointer-generator, qui combinent generation et copie depuis la source. Cette idee est importante pour les resumes d'actualites, car les noms propres, les lieux, les chiffres et les evenements doivent souvent etre repris fidelement depuis le texte original \cite{See2017PointerGenerator}. Ces travaux montrent deja une tension centrale pour notre projet: un bon systeme de resume doit etre capable de reformuler, mais il doit aussi rester proche de la source pour eviter les erreurs factuelles.

Dans le contexte du present projet, cette evolution justifie le choix de modeles de generation modernes. Le but n'est pas seulement de raccourcir un article, mais de produire un resume lisible, informatif et utilisable dans une application web. Les approches extractives simples peuvent etre utiles comme solution de secours, mais elles ne suffisent pas pour obtenir des resumes fluides dans plusieurs langues. C'est pourquoi le projet s'appuie principalement sur des modeles transformer sequence-a-sequence, tout en gardant des mecanismes de controle comme le nettoyage des textes, le post-traitement et l'evaluation de la fidelite a la source.

\section{Resume d'actualites et specificites du domaine}

Le resume d'actualites est un cas particulier du resume automatique. Les articles de presse contiennent souvent des informations factuelles, des dates, des noms de personnes, des lieux, des citations et des chiffres. Une erreur dans un resume d'actualite peut donc modifier le sens de l'information. Les jeux de donnees comme CNN/Daily Mail, XSum, Newsroom et Multi-News ont joue un role important dans le developpement des modeles de resume pour le domaine journalistique \cite{Hermann2015Teaching,Narayan2018XSum,Grusky2018Newsroom,Fabbri2019MultiNews}. CNN/Daily Mail a ete largement utilise pour des resumes relativement extractifs, alors que XSum se concentre sur des resumes courts et plus abstractive, souvent proches d'une phrase de synthese \cite{Narayan2018XSum}.

Le jeu de donnees Newsroom est particulierement utile pour comprendre que les resumes d'actualites ne sont pas tous du meme type. Grusky et al. montrent que certains resumes sont tres proches du texte source, tandis que d'autres reformulent davantage l'information \cite{Grusky2018Newsroom}. Cette observation est directement liee a notre projet: les modeles peuvent avoir des comportements differents selon leur degre d'extractivite, leur capacite a reformuler et leur fidelite au texte original. Un modele tres extractif peut etre plus prudent, mais moins naturel; un modele plus abstractive peut etre plus fluide, mais aussi plus risque du point de vue factuel.

Ces travaux montrent aussi une limite importante de la litterature: beaucoup d'etudes evaluent les modeles sur des jeux de donnees fixes et propres, alors que les textes provenant de flux RSS et de sites d'actualites reels sont plus variables. Ils peuvent contenir des titres repetes, des informations de navigation, des elements publicitaires, des phrases incompletes ou des structures HTML complexes. Le present projet cherche a reduire cet ecart entre evaluation academique et usage pratique en integrant les modeles dans une application qui collecte de vrais articles, extrait le corps du texte, nettoie les contenus et produit des resumes utilisables par l'utilisateur final.

\section{Modeles transformer pour le resume}

Les modeles transformer ont profondement transforme les taches de generation de texte. Contrairement aux architectures recurrentes, ils gerent mieux les dependances longues et peuvent etre pre-entraines sur de grands corpus avant d'etre adaptes a une tache specifique. BART est l'un des modeles les plus importants pour le resume abstractive. Il est pre-entraine comme un autoencodeur denoisant: le modele apprend a reconstruire un texte original a partir d'une version degradee, ce qui le rend efficace pour des taches de generation comme le resume \cite{Lewis2020BART}. Cette logique explique pourquoi les modeles BART sont souvent utilises comme base pour le resume d'articles en anglais.

PEGASUS propose une autre strategie de pre-entrainement, plus directement orientee vers le resume. Le modele apprend a predire des phrases importantes supprimees d'un document, ce qui rapproche le pre-entrainement de la tache finale de summarisation \cite{Zhang2020PEGASUS}. Meme si PEGASUS n'est pas le modele central de notre implementation, cette etude montre que la qualite d'un modele de resume depend non seulement de son architecture, mais aussi de la maniere dont il est pre-entraine et adapte a la tache.

Dans notre projet, les modeles BART servent surtout de reference pour le resume en anglais, tandis que l'objectif principal est le resume multilingue. Cette distinction est importante: un modele performant en anglais ne peut pas etre directement considere comme une solution generale pour des langues comme le turc, l'arabe, le russe, le chinois ou le vietnamien. Pour couvrir plusieurs langues dans une meme application, il faut utiliser des modeles pre-entraines sur des donnees multilingues et capables de generer dans differentes langues.

\section{Modeles multilingues et resume multilingue}

Les modeles multilingues comme mBART, mBART-50 et mT5 sont centraux pour le present projet. mBART etend l'idee du pre-entrainement denoisant a plusieurs langues. Le modele apprend des representations partagees qui peuvent ensuite etre adaptees a des taches de generation multilingues \cite{Liu2020mBART}. mBART-50 augmente encore cette capacite en couvrant davantage de langues et en facilitant le transfert multilingue \cite{Tang2020mBART50}. Ces caracteristiques rendent mBART adapte au resume d'actualites multilingue, car un meme cadre de modele peut etre utilise pour plusieurs langues.

mT5 suit une autre approche: il adapte le modele text-to-text T5 a un contexte massivement multilingue \cite{Xue2021mT5}. Dans notre projet, le modele mT5 utilise pour la comparaison n'a pas ete entraine dans le cadre du projet. Il a ete pris comme modele deja entraine dans une etude similaire et sert de reference externe. Cette distinction est importante pour l'interpretation des resultats: mT5 permet de situer les modeles mBART entraines dans le projet face a un modele multilingue fort, mais il ne doit pas etre presente comme une sortie de notre propre processus d'entrainement.

Les jeux de donnees multilingues jouent aussi un role essentiel. XL-Sum est l'un des principaux jeux de donnees pour le resume abstractive multilingue d'actualites. Il couvre un grand nombre de langues et fournit des couples article-resume adaptes a l'entrainement et a l'evaluation de modeles multilingues \cite{Hasan2021XLSum}. MLSUM est un autre corpus multilingue important, construit a partir d'articles de presse dans plusieurs langues europeennes \cite{Scialom2020MLSUM}. Ces ressources montrent que le resume multilingue est devenu un sujet important, mais elles restent principalement utilisees dans des evaluations hors ligne. Notre projet complete cette perspective en reliant l'evaluation multilingue a une application web qui traite des flux d'actualites reels.

\section{Evaluation automatique des resumes}

L'evaluation est l'un des problemes les plus importants dans le resume automatique. Les metriques classiques comme ROUGE, BLEU et METEOR sont tres utilisees parce qu'elles permettent une comparaison automatique et reproductible \cite{Lin2004ROUGE,Papineni2002BLEU,Banerjee2005METEOR}. ROUGE mesure principalement le recouvrement de n-grammes entre le resume produit et le resume de reference. BLEU, propose initialement pour la traduction automatique, mesure aussi la correspondance de n-grammes. METEOR ajoute une forme plus flexible de correspondance entre tokens et peut mieux tenir compte de certaines variations lexicales.

Cependant, ces metriques ont des limites importantes pour le resume d'actualites. Un resume peut obtenir un bon score ROUGE parce qu'il partage des mots avec la reference, mais il peut quand meme etre incomplet, repetitif ou peu fidele au texte source. A l'inverse, un resume bien formule peut utiliser des mots differents de la reference et recevoir un score plus faible. SummEval montre que l'evaluation du resume doit prendre en compte plusieurs dimensions, notamment la coherence, la pertinence, la fluidite et la coherence factuelle \cite{Fabbri2021SummEval}. Cette observation soutient directement la methode adoptee dans notre projet.

Une etude recente de Miller et Tang renforce aussi cette critique des evaluations classiques. Les auteurs soutiennent que les metriques appliquees aux modeles generatifs doivent etre alignees avec les usages reels plutot qu'avec des benchmarks abstraits \cite{Miller2025LLMMetrics}. Leur analyse propose d'examiner les capacites des modeles selon des criteres pratiques comme la coherence, l'exactitude, la clarte, la pertinence et l'efficacite. Cette perspective est directement liee au present projet, car les scores proxy utilises dans l'evaluation des modeles de resume cherchent egalement a mesurer des dimensions proches de l'usage reel, au-dela du simple recouvrement lexical.

Pour cette raison, le projet ne se limite pas aux metriques classiques. Il introduit des scores proxy composites pour mesurer des dimensions pratiques de la qualite: coherence, exactitude, clarte, pertinence, efficacite, fidelite a la source et completude du contenu. Ces scores sont construits a partir de signaux mesurables comme la similarite avec la reference, la couverture de la source, le rappel de la source, la repetition, la longueur du resume, le taux de compression et le temps de generation. Ils ne remplacent pas une evaluation humaine, mais ils permettent une comparaison automatique plus riche entre les modeles. Ce choix repond a un manque dans les evaluations classiques: mesurer non seulement la proximite avec une reference, mais aussi l'utilisabilite du resume dans une application reelle.

\section{Fidelite factuelle et hallucination dans les resumes}

La fidelite factuelle est un enjeu majeur pour le resume abstractive. Les modeles generatifs peuvent produire des informations qui ne sont pas presentes dans le texte source, modifier des relations causales, confondre des entites ou changer des chiffres. Maynez et al. montrent que les resumes abstractive peuvent contenir des erreurs factuelles meme lorsqu'ils semblent bien formes linguistiquement \cite{Maynez2020Faithfulness}. Ce probleme est particulierement sensible pour les actualites, car l'utilisateur peut prendre le resume comme une representation rapide d'un evenement reel.

Plusieurs travaux ont propose des methodes pour mesurer ou detecter ces inconsistances. FactCC utilise une approche de classification pour evaluer la coherence factuelle entre un document et son resume \cite{Kryscinski2020FactCC}. SummaC s'appuie sur des modeles d'inference en langue naturelle pour detecter les contradictions ou les inconsistances \cite{Laban2022SummaC}. QuestEval propose une evaluation basee sur des questions et reponses afin d'examiner si les informations importantes du texte source sont bien preservees \cite{Scialom2021QuestEval}. Ces travaux montrent que la fidelite ne peut pas etre reduite a un simple recouvrement lexical.

Le present projet ne construit pas un systeme complet de fact-checking, mais il prend en compte ce probleme dans sa methode. Le score \textit{quality factuality}, la couverture de la source et les mesures de fragments extraits servent de proxy pour estimer si le resume reste proche du texte original. De plus, l'application conserve le lien vers l'article source afin que l'utilisateur puisse verifier le contexte complet. Cette decision est importante sur le plan methodologique et ethique: le resume est presente comme une aide a la lecture rapide, non comme un remplacement definitif de l'article.

\section{Systemes de resume d'actualites et integration applicative}

La litterature contient aussi des travaux sur les systemes de resume d'actualites. Newsblaster est un exemple ancien mais important: il collectait automatiquement des nouvelles, les groupait par sujet et produisait des resumes pour aider les utilisateurs a suivre l'actualite \cite{McKeown2002Newsblaster}. MEAD a egalement propose une plateforme pour le resume multi-document et multilingue, avec une attention particuliere aux systemes configurables \cite{Radev2004MEAD}. Ces systemes montrent que le resume automatique n'est pas seulement un probleme de modele, mais aussi un probleme d'architecture, de collecte, de filtrage et de presentation des resultats.

Notre projet s'inscrit dans cette tradition applicative, mais avec des modeles neuronaux modernes et une architecture orientee flux RSS. L'application developpee collecte les articles a partir de sources definies, extrait le texte principal, nettoie les donnees, genere des resumes avec plusieurs modeles et stocke les resultats dans une base SQLite. L'utilisation d'un processus d'ingestion en arriere-plan evite de faire attendre l'utilisateur pendant la generation du resume. Cette approche repond a une limite pratique souvent absente des evaluations academiques: un modele doit non seulement produire de bons scores, mais aussi fonctionner dans un systeme interactif avec des contraintes de temps, de stockage et de qualite des entrees.

\section{Lacunes identifiees et positionnement du projet}

L'analyse de la litterature montre plusieurs lacunes. Premierement, beaucoup de travaux evaluent les modeles sur des corpus propres et fixes, alors que les textes d'actualites reels sont plus instables et plus bruites. Deuxiemement, les evaluations reposent souvent sur des metriques classiques comme ROUGE ou BLEU, qui ne capturent pas suffisamment la fidelite a la source, la lisibilite, la repetition, la vitesse ou l'adaptation a une application reelle. Troisiemement, les travaux sur le resume multilingue se concentrent souvent sur la performance des modeles, tandis que l'integration dans une application web de flux d'actualites est moins souvent analysee.

Ce projet cherche a combler ces lacunes en combinant quatre elements. D'abord, des modeles mBART multilingues sont entraines et adaptes au resume d'actualites. Ensuite, un modele mT5 deja entraine dans une etude similaire est utilise comme reference externe pour la comparaison. Puis, une evaluation multidimensionnelle est proposee avec des metriques classiques et de nouveaux scores proxy. Enfin, les modeles sont integres dans une application web capable de collecter, resumer, stocker et afficher des actualites issues de flux RSS. Cette combinaison donne au projet une valeur a la fois scientifique et technologique: il analyse les modeles de maniere experimentale, mais il verifie aussi leur utilite dans un systeme logiciel fonctionnel.

Ainsi, la revue de litterature montre que le sujet du projet est important parce qu'il se situe a l'intersection de trois besoins actuels: resumer rapidement de grands volumes d'actualites, traiter plusieurs langues dans un meme systeme et evaluer les modeles selon des criteres plus proches de l'usage reel. Le projet ne remplace pas les approches existantes; il les prolonge en reliant l'entrainement de modeles, l'evaluation critique et l'integration applicative.
